% part: intuitionistic-logic
% chapter: propositions-as-types
% section: introduction

\documentclass[../../../include/open-logic-section]{subfiles}

\begin{document}

\olfileid{int}{pty}{int}

\olsection{Introduction}

The lambda calculus is a simple model of computation. It operates on
terms built from variables. If $N$ and $M$ are terms, then $(NM)$ is
also a term (``$N$ applied to~$M$''). If $N$ is a term, then
$\lambd[x][N]$ is a term. If $N$ contains the variable~$x$, then the
corresponding occurrences of~$x$ in $\lambd[x][N]$ are bound by the
$\lambd[x]$ operator, and so are no longer free. A term of the form
$(\lambd[x][N]M)$ is said to $\beta$-contract to $\Subst{N}{M}{x}$
(the result of replacing every free occurrence of $x$ in~$N$ by $M$).
A term $N$ $\beta$-converts to~$N'$ if $N'$ results from $N$ by
$\beta$-contracting a subterm, and we write $N \redone N'$. A
computation in the $\lambd$-calculus is a sequence $N \redone N_1
\redone N_2 \redone \dots$. If this sequence ends in a term $N_k$
where no subterm can be $\beta$-contracted, it is called \emph{normal}
or a normal form of~$N$. Think of computation in the $\lambd$-calculus
as such sequences. The computation halts if it results in a normal
form, and the normal form is the result of the computation. This
simple model of computations turns out to be Turing complete: every
partial recursive function can be computated by a term in the $\lambd$-calculus.

Haskell Curry and William Howard independently discovered a close
similarity between natural deduction and the $\lambd$-calculus.
Intuitively, a term $\lambd[x][N]$ is a function: $x$ is the argument;
$N$ tells us how to compute the value given~$x$. A term
$(\lambd[x][N]M)$ then is applying the function $\lambd[x][N]$ to an
argument~$M$, and $\beta$-contraction tells us how to evaluate it:
replace $x$ in $N$ by~$M$ (and keep evaluating the resulting term).
Now think of these as functions with a fixed domain and range. If the
domain of $\lambd[x][N]$ is $A$ and and its range~$B$, then we should
think of the variable~$x$ as only taking arguments from~$A$ and $N$ as
producing values from~$B$. Consequently, $(\lambd[x][N]M)$ only should
make sense if $\lambd[x][N]$ is a function $A \to B$, and $M \in A$.
If we want to add this information to the $\lambd$-calculus we obtain
the \emph{typed} $\lambd$-calculus. In the typed $\lambd$-calculus not
every expression $(\lambd[x][N]M)$ is legal, but only those where the
``types'' of $\lambd[x][N]$ and $M$ fit together, i.e., $\lambd[x][N]$
has type $A \lif B$ and $M$ has type~$A$. In fact, not every term
$(M_1M_2)$ is legal only if $x$ is of type $A$ and $N$ is of type only
those where $M_1$ has type $A \to B$ and $M_2$ has type~$A$; if so,
then $(M_1M_2)$ has type~$B$. And if the type of $x$ is $A$ and the
type of $N$ is $B$ then the type of $\lambd[x][N]$ is $A \to B$.

These typing rules now clearly correspond to !!{derivation}s in
natural deduction, where the types are represented by !!{formula}s,
and the !!{derivation}s themselves to $\lambd$-terms. For instance,
!!a{derivation} of $!A \lif !B$ corresponds to a term
$\lambd[\typeof{x}{!A}][N]$, which has the function type
$!A \lif !B$. The inference rules of natural deduction correspond to
typing rules in the typed lambda calculus, e.g.,
\begin{prooftree}
  \AxiomC{$\Discharge{!A}{x}$}
  \DeduceC{$!B$}
  \DischargeRule{\Intro\lif}{x}
  \UnaryInfC{$!A \lif !B$}
  \DisplayProof\bottomAlignProof
  \quad\text{corresponds to}\quad
  \Axiom$x: !A \fCenter N: !B$
  \UnaryInf$ \fCenter \lambd[\typeof{x}{!A}][N] : !A \lif !B$
\end{prooftree}
where the rule on the right means that if $x$ is of type~$!A$ and $N$
is of type~$!B$, then $\lambd[\typeof{x}{!A}][N]$ is of type~$!A \lif
!B$.

The $\Elim\lif$ rule corresponds to the typing rule for composition
terms, i.e.,
\[
  \AxiomC{$!A \lif !B$}
  \AxiomC{$!A$}
  \RightLabel{\Elim\lif}
  \BinaryInfC{$!B$}
  \DisplayProof
  \quad\text{corresponds to}\quad
  \Axiom$\fCenter M_1: !A \lif !B$
  \Axiom$\fCenter M_2: !A$
  \BinaryInf$ \fCenter (M_1M_2) : !B$
  \DisplayProof
\]
If a \Intro\lif{} rule is followed immediately by a \Elim\lif{} rule,
the !!{derivation} can be simplified:
\begin{gather*}
  \AxiomC{$\Discharge{!A}{x}$}
  \DeduceC{$!B$}
  \DischargeRule{\Intro\lif}{x}
  \UnaryInfC{$!A \lif !B$}
  \AxiomC{}
  \DeduceC{$!A$}
  \RightLabel{\Elim\lif}
  \BinaryInfC{$!B$}
  \DisplayProof
  \quad\redone\quad
  \AxiomC{}
  \DeduceC{$!A$}
  \DeduceC{$!B$}
  \DisplayProof
\end{gather*}
which corresponds to $\beta$-contraction of lambda terms
\[
(\lambd[\typeof{x}{!A}][N]M) \quad\redone\quad
\Subst{N}{M}{x}.
\]

Similar correspondences hold between the rules for $\land$ and ``product''
types, and between the rules for $\lor$ and ``sum'' types.

This correspondence between terms in the simply typed lambda calculus
and natural deduction !!{derivation}s is called the ``Curry--Howard'',
``proofs-as-programs,'' or ``propositions as types'' correspondence.
In addition to !!{formula}s (propositions) corresponding to types, and
proofs to terms, we can summarize the correspondences as follows:
\begin{center}
  \begin{tabular}{c c c}
    logic & program \\
    \hline
    proposition/!!{formula} & type \\
    proof & term \\
    assumption & variable \\
    discharged assumption & bond variable \\
    undischarged assumption & free variable \\
    $!A \lif !B$ & function type \\
    $!A \land !B$ & product type \\
    $!A \lor !B$ & sum type \\
    $\lfalse$ & empty type \\
    \hline
  \end{tabular}
\end{center}

The Curry--Howard correspondence is one of the cornerstones of
automated proof assistants and type checkers for programs, since
checking a proof witnessing a proposition (as we did above) amounts to
checking if a program (term) has the declared type.
\end{document}
